\documentclass[11pt]{article}

% ---------------------------------------------------------------------------
%  L3 — Bounds for the polyplet growth constant.
%  Category L (entirely machine-written); see paper/README.md and
%  docs/publication-split.md.
%
%  Source material, all in this repository:
%    docs/proofs/polyplet-upper-bound.md       the whole upper-bound story
%    results/growth-constant.md          the Collatz-Wielandt certificates
%    results/growth-constant.md     the mu_H ladder and the estimate
%    results/strip_mu_certificates.log         the receipts quoted in Table 2
%    results/growth-constant.md      the closed alternative route
%    results/growth-constant.md       what is still open
%    docs/lean-record.md                     what Lean does and does not cover
% ---------------------------------------------------------------------------

% ---------------------------------------------------------------------------
%  paper/shared/preamble.tex — packages and macros common to every manuscript
%  in this directory.  Factored out of paper/polyplets-report.tex and
%  paper/technical-report.tex, which between them had two copies of the same
%  twenty lines.
%
%  technical-report.tex does NOT \input this file and is not to be edited to do
%  so: it is jasonp's prose and is read-only to the machine (paper/README.md,
%  "Who may edit what").  The duplication there is deliberate.
%
%  Assumes \documentclass[11pt]{article} has already been issued.
% ---------------------------------------------------------------------------
\usepackage[margin=1in]{geometry}
\usepackage{amsmath,amssymb}
\usepackage{amsthm}
\usepackage{booktabs}
\usepackage{graphicx}
\usepackage{numprint}
\usepackage{microtype}
\usepackage[table]{xcolor}
\usepackage[hidelinks]{hyperref}
\usepackage{flafter}   % a float never appears before the text that introduces it
\usepackage{float}     % [H] pins tables/figures exactly where written
\usepackage[font=small,labelfont=bf,labelsep=period]{caption}
\setlength{\intextsep}{16pt plus 3pt minus 3pt}

\npthousandsep{,}
\npfourdigitsep
\newcommand{\num}[1]{\numprint{#1}}
\newcommand{\g}{\cellcolor{gray!15}}

% Theorem environments.  One counter shared across kinds, so that a reference
% like "Theorem 4" is unambiguous in a paper that also has lemmas and
% propositions.
\theoremstyle{plain}
\newtheorem{theorem}{Theorem}
\newtheorem{proposition}[theorem]{Proposition}
\newtheorem{lemma}[theorem]{Lemma}
\newtheorem{corollary}[theorem]{Corollary}
\theoremstyle{definition}
\newtheorem{definition}[theorem]{Definition}
\newtheorem{example}[theorem]{Example}
\theoremstyle{remark}
\newtheorem{remark}[theorem]{Remark}
\newtheorem{conjecture}[theorem]{Conjecture}
\newtheorem{openproblem}[theorem]{Open Problem}

% Repo-path citations.  Every one of these must resolve in the tree that ships
% with the paper; `make gate-citations` checks the markdown, and
% paper/verify_claims.py checks the manuscripts.
\newcommand{\repo}[1]{\texttt{#1}}

% OEIS links.
\newcommand{\oeis}[1]{\href{https://oeis.org/#1}{#1}}

% Growth constants, so a rename is one edit.
\newcommand{\lam}{\lambda}
\newcommand{\muH}{\mu_H}

% ---------------------------------------------------------------------------
%  paper/shared/disclosure.tex — the two authorship disclosure blocks.
%
%  docs/publication-split.md §1 fixes the rule these implement:
%
%    "Under no circumstances will there be any ambiguity about whether a paper
%     is wholly my work (none), merely my prose (some) or entirely LLM-authored
%     (some)."   — jasonp, 2026-08-07
%
%  There are exactly two blocks and they live in exactly one file, so that a
%  paper cannot quietly soften its own disclosure.  The fixed sentences take no
%  arguments.  The ONE thing a paper supplies is the per-result verification
%  ledger, because that is the part that differs honestly between papers — and
%  §1 requires it to be per result, not in aggregate.
%
%  Do not add a \Pdisclosure or \Ldisclosure variant with an optional softening
%  argument.  If a paper does not fit one of these two blocks, the paper is
%  wrong, not the block.
% ---------------------------------------------------------------------------

\newsavebox{\disclosurebox}

\newenvironment{disclosureframe}
  {\begin{center}\begin{lrbox}{\disclosurebox}\begin{minipage}{0.93\textwidth}\small}
  {\end{minipage}\end{lrbox}\fbox{\usebox{\disclosurebox}}\end{center}}

% --- P: jasonp's prose -------------------------------------------------------
% Byline: Jason Parker, alone.  Argument: the per-result ledger of what the
% machine did.
\newcommand{\Pdisclosure}[1]{%
\begin{disclosureframe}
\textbf{Authorship.}\enspace Every sentence of this paper was written by its
named author. He has read every proof it states and believes each one, and he
is responsible for the correctness of what it claims.

\smallskip
\textbf{Machine contribution.}\enspace #1
\end{disclosureframe}}

% --- L: entirely machine-written --------------------------------------------
% Argument: the per-result verification ledger — for each result, what warrants
% it (Lean kernel, exact certificate, or labelled measurement) and how much of
% it a human has checked.
\newcommand{\Ldisclosure}[1]{%
\begin{disclosureframe}
\textbf{Authorship.}\enspace This document was written end to end by a large
language model. That includes its mathematics, its prose, its tables and its
\LaTeX{}. No sentence of it was written by a human. The mathematics it reports
was likewise derived by machine.

\smallskip
\textbf{Human role.}\enspace The person who directed this work chose what to
compute, chose what to publish, and is responsible for the decision to release
this document. Except where the ledger below says otherwise, that person has
not personally checked the arguments here, and this disclosure is not to be
read as an endorsement of them.

\smallskip
\textbf{What warrants each result.}\enspace #1
\end{disclosureframe}}

% --- Draft banner ------------------------------------------------------------
% For an L-paper whose verification ledger is not yet settled.  Loud on purpose:
% an L-paper without a truthful ledger must not be mistaken for a finished one.
%
% \firstdraftbanner is the standard wording, which five papers previously
% carried character-for-character in their own preambles.  It lives here for the
% same reason the disclosure blocks do: identical text in seven places drifts,
% and a paper must not be able to soften it by editing its own copy.  The
% optional argument is for a gate specific to one paper, appended verbatim; a
% paper needing different wording (L6, compute-gated) calls \draftbanner direct.
\newcommand{\firstdraftbanner}[1][]{%
\draftbanner{This is a first draft. The verification ledger below records what
warrants each result \emph{mechanically}; the column that records what a human
has read is empty, deliberately, because nobody has read it yet. Do not
circulate until that column says something true.#1}}

\newcommand{\draftbanner}[1]{%
\begin{center}
\fcolorbox{red!60!black}{yellow!15}{%
  \begin{minipage}{0.93\textwidth}\small
  \textbf{DRAFT — NOT FOR CIRCULATION.}\enspace #1
  \end{minipage}}
\end{center}}


\newcommand{\Z}{\mathbb{Z}}
\newcommand{\Q}{\mathbb{Q}}
\newcommand{\R}{\mathbb{R}}
\newcommand{\rhoo}{\rho}

\title{\textbf{A certified two-sided bound for the polyplet growth constant}\\[4pt]
\large $6.543 \le \lambda \le 9.3154$, both ends checkable in exact
integer arithmetic}
\author{[byline: jasonp's call --- \repo{docs/publication-split.md} \S1]}
\date{August 2026}

\begin{document}
\maketitle

\firstdraftbanner

\Ldisclosure{%
\emph{Existence of $\lambda$} --- Fekete's lemma applied to a
supermultiplicativity that is proved in Lean~4 against mathlib
(\texttt{a\_supermul}, \texttt{lambda\_tendsto}, \texttt{a\_le\_lambda\_pow},
\texttt{lambda\_le}), standard axioms only, no \texttt{sorry}.
\emph{The lower bound $\lambda \ge 6.543$} --- an exact rational certificate:
a nonnegative integer vector and a rational $x$ for which $M(x)v \ge v$ is
verified componentwise in exact $128$-bit integer arithmetic, with every
rounding downward. The receipt is
\repo{results/strip\_mu\_certificates.log}; the checker carries a four-part
self-test that includes a case it must reject
(\S\ref{sec:selftest}). The transfer operator itself is the production kernel
and is \emph{trusted, not certified} --- see \S\ref{sec:honest}.
\emph{The upper bound $\lambda \le 9.3154$} --- an exact rational
super-solution of a $5930$-component convolution system, all inequalities
checked in exact rational arithmetic with no floating point anywhere in the
chain; formalized in Lean (\texttt{RatCert.lambda\_le},
\texttt{lambda\_le\_of\_buiSystem}: standard axioms only; the concrete
certificate \texttt{lambda\_le\_of\_bui\_rd3} adds one named
\texttt{native\_decide} leaf, which is the certificate's arithmetic validity
check).
\emph{The elementary bound $\lambda \le 5^5/4^4$
(Proposition~\ref{prop:crude})} --- proved, by a BFS-frame encoding that is
machine-verified over every animal with $n \le 8$ with a negative control. Its
\emph{earlier} proof, by the Redelmeier decision string, was found to be false
on re-entrant animals and is retained in \S\ref{sec:crudebound} as a warning
rather than deleted; the constant never moved.
\emph{The slack audit of \S\ref{sec:gap}} --- labeled measurement against
brute-force counts, not a theorem.
\emph{The estimate $\lambda \approx 7.11$} --- labeled measurement
(extrapolation), and used nowhere in the proofs.
\emph{Human verification:} none, as of this draft.}

\begin{abstract}
\noindent
A \emph{polyplet}, or king animal, is a finite set of cells of the square
lattice connected through shared edges or corners; $a(n)$ counts those with $n$
cells up to translation (OEIS \oeis{A006770}). The limit
$\lambda = \lim_n a(n)^{1/n}$ exists by supermultiplicativity. We prove
\[
  6.543 \;\le\; \lambda \;\le\; 9.3154 ,
\]
with both ends reduced to finite integer computations that a reader can re-run
without repeating any of the reasoning that produced them.

The lower end is a ladder. The height-$H$ strips form a subfamily, so their
growth constant $\mu_H$ bounds $\lambda$ from below; we certify $\mu_H$ for
$H \le 17$ by exhibiting a nonnegative integer vector $v$ and a rational $x$
with $M(x)v \ge v$ componentwise, checked in exact $128$-bit arithmetic with
every rounding taken downward. This is Collatz--Wielandt in its lower half only,
which needs neither irreducibility nor positivity. The top rung,
$\mu_{17} \ge 6.543$, beats the best previously published lower bound, with a
certificate where that value was numerical.

We found no published upper bound to compare against. Elementary
counting gives $\lambda \le 5^5/4^4 \approx 12.2$ and, we show, cannot be pushed
below that: the two obvious improvements are both non-injective, and both are
caught by the same control, namely that the identical argument would prove
$\lambda_{\text{polyomino}} \le 4$, which is false. We instead port Bui's
convolution-certificate method to king adjacency and exhibit an exact rational
super-solution of the resulting $5930$-component system, giving
$\lambda \le 9.3154$ in exact arithmetic.

We also report, as measurement rather than theorem, why the upper bound stops
where it does. An audit of every recurrence in the system shows the over-count
is diffuse (median slack $1.144$, maximum $1.222$) and that it
\emph{grows with $n$} --- so it is not a local defect that a larger window or
per-type tuning could remove, but the same eight-neighbor connectivity
obstruction that closed the other routes we tried.
\end{abstract}

\tableofcontents

\section{Introduction}

Let $a(n)$ be the number of fixed polyplets with $n$ cells: finite sets of unit
cells of the square lattice, connected through shared edges \emph{or corners},
counted up to translation. This is OEIS \oeis{A006770}, and it begins
$1, 4, 20, 110, 638, 3832, 23592, \dots$

Concatenating an animal of $m$ cells with one of $n$ cells in a fixed relative
position, in a way that keeps the result connected and keeps the map injective,
gives $a(m)a(n) \le a(m+n)$; Fekete's lemma then makes
$\lambda := \lim_n a(n)^{1/n}$ exist, with $a(n) \le \lambda^n$ for every $n$ ---
so every enumerated term is a floor on $\lambda$ and never an approximation to
it. The largest such floor is $a(40)^{1/40} = 6.2208$ and the last successive
ratio is $a(40)/a(39) = 6.9352$; extrapolating that ratio sequence, by
differential approximants on the 40 terms, gives the estimate
$\lambda = 7.110(1)$, which is a measurement and not a bound.

The state of rigorous knowledge before this work was one-sided. Bacher's directed
animals on the king lattice give $\lambda \ge 3 + 2\sqrt2 \approx 5.828$ in closed
form~\cite{bacher2015}, and his multi-directed animals give $\approx 6.4752$,
that one numerically rather than in closed form. No upper bound on $\lambda$ for
the king lattice appears in the literature we searched; the searches are
recorded in \repo{literature/MISSING.md}. We claim the bound, not priority.

\subsection{What is proved}

\begin{theorem}
\label{thm:main}
$\dfrac{6543}{1000} \;\le\; \lambda \;\le\; \dfrac{20000}{2147}$, that is
$6.543 \le \lambda \le 9.3154$.
\end{theorem}

Both decimals round in the safe direction, and the rationals are what is
actually certified.\footnote{The lower certificate proves
$\lambda \ge 6543/1000$, so its decimal is exact. The upper certificate proves
$\lambda \le 20000/2147 = 9.31532\dots$, so the four-decimal form of that bound
is $9.3154$ and not $9.3153$; truncating an upper bound would claim
$3.2\times10^{-5}$ more than the certificate gives.}

Each half is a \emph{certificate}: a finite object, an integer vector on one
side and a rational vector on the other, together with a finite integer or
rational computation checking a system of inequalities against it. The searches
that found those objects were floating-point and are not part of the proof. A
worse object would still check; a wrong one fails the check and names the
failing component.

\section{Preliminaries}

\begin{definition}
A \emph{polyplet} of size $n$ is a connected set of $n$ cells of $\Z^2$ under
king adjacency (the eight offsets $(\pm1,0),(0,\pm1),(\pm1,\pm1)$), taken up to
translation. $a(n)$ is their number; $\lambda := \lim_n a(n)^{1/n}$.
\end{definition}

Supermultiplicativity, the existence of the limit, the per-term inequality and
the resulting lower bound are formalized in Lean~4 against
mathlib~\cite{mathlib} and depend on standard axioms only. The formalization
carries them once, generic in the sequence, so that $\lambda$ and the growth
constant of a staircase subfamily are two instances of the same development.

\begin{definition}[strip counts and strip growth constants]
\label{def:strip}
$T_{\le H}(n)$ is the number of polyplets of size $n$ whose bounding box has
height at most $H$, and
$\mu_H := \lim_n T_{\le H}(n)^{1/n}$.
\end{definition}

Each $\mu_H$ is the reciprocal of the radius of convergence of a rational
generating function, the height-$\le H$ column-signature transfer matrix being
finite, so it is an algebraic number, and the sequence is strictly increasing
and bounded by $\lambda$.

\section{The lower bound}
\label{sec:lower}

\subsection{The three inequalities}

The lower bound is the composition of three facts. The first two are about
nonnegative matrices; only the third is about polyplets.

\paragraph{(1) Collatz--Wielandt, lower half.} Let $A$ be a nonnegative square
matrix and let $v \ge 0$, $v \ne 0$, satisfy $Av \ge v$ componentwise. Then
$\rhoo(A) \ge 1$. Indeed $A$ is monotone on the nonnegative orthant, so
$A^k v \ge v$ for every $k$; picking any $i$ with $v_i > 0$ gives
$\|A^k\|_\infty \ge v_i/\|v\|_\infty$, and $\rhoo(A) = \lim_k \|A^k\|^{1/k} \ge 1$.

This requires \emph{neither} positivity of $v$ nor
irreducibility of $A$. Only the lower half of the two-sided Collatz--Wielandt
characterisation is used; the upper half, which does need irreducibility, never
appears. See~\cite[Ch.~8]{hornJohnson2013} or~\cite[Ch.~2]{bermanPlemmons1994}.

\paragraph{(2) Monotonicity in $x$.} Let $M(x)$ be the height-$\le H$ king
connectivity transfer operator with entries $x^{\text{cells placed}}$: entrywise
nonnegative and entrywise increasing in $x > 0$. The spectral radius of a
nonnegative matrix is monotone in its entries, so $\rhoo(M(x))$ increases in $x$,
and the point $x^*$ where $\rhoo = 1$ is the radius of convergence of the
height-$\le H$ generating function, with $\mu_H = 1/x^*$. Therefore
\[
  \rhoo(M(x)) \ge 1 \;\Longrightarrow\; x \ge x^*
             \;\Longrightarrow\; \mu_H = 1/x^* \ge 1/x .
\]

\paragraph{(3) Subfamily containment.} Height-$\le H$ polyplets are polyplets, so
$T_{\le H}(n) \le a(n)$ for every $n$ and hence $\mu_H \le \lambda$.

Composing: an exactly checked inequality $M(\mathit{den}/\mathit{num})\,v \ge v$
certifies $\mu_H \ge \mathit{num}/\mathit{den}$, and therefore
$\lambda \ge \mathit{num}/\mathit{den}$.

\begin{definition}[certificate]
\label{def:cert}
A \emph{certificate} for the rational $r = \mathit{num}/\mathit{den}$ at height
$H$ is a nonnegative, nonzero integer vector $v$ indexed by the states of
$M$, together with the verified statement that
$M(\mathit{den}/\mathit{num})\,v \ge v$ holds componentwise in exact integer
arithmetic.
\end{definition}

\subsection{The certified ladder}

\begin{theorem}
\label{thm:ladder}
The values of Table~\ref{tab:ladder} are certified in the sense of
Definition~\ref{def:cert}. In particular $\mu_{17} \ge 6.543$, and hence
$\lambda \ge 6.543$.
\end{theorem}

\begin{table}[H]
\centering\small
\begin{tabular}{r r r r r r}
\toprule
$H$ & certified $\mu_H \ge$ & states & $\mu_H$ (float) & tries & wall (s) \\
\midrule
 2 & 2.4142135 &         3 & 2.414213562 &  1 &   0.0 \\
 3 & 3.4437183 &         8 & 3.443718375 &  1 &   0.0 \\
 4 & 4.1823214 &        20 & 4.182321413 &  1 &   0.0 \\
 5 & 4.7178012 &        50 & 4.717801291 &  1 &   0.0 \\
 6 & 5.1153244 &       126 & 5.115324460 &  1 &   0.0 \\
 7 & 5.4178476 &       322 & 5.417847610 &  1 &   0.0 \\
 8 & 5.6533727 &       834 & 5.653372762 &  1 &   0.0 \\
 9 & 5.8404579 &      2187 & 5.840457941 &  1 &   0.1 \\
10 & 5.9916957 &      5797 & 5.991695792 &  1 &   0.2 \\
11 & 6.1158416 &     15510 & 6.115841628 &  1 &   0.5 \\
12 & 6.2191246 &     41834 & 6.219124621 &  1 &   1.7 \\
13 & 6.3060712 &    113633 & 6.306071285 &  1 &   5.4 \\
14 & 6.3800149 &    310571 & 6.380034445 & 18 &  18.3 \\
15 & 6.443532  &    853466 & 6.443540832 &  8 &  57.7 \\
16 & 6.4984    &   2356778 & 6.498524533 &  3 & 225.0 \\
17 & \textbf{6.543} & 6536381 & 6.546486999 &  6 & 772.7 \\
\bottomrule
\end{tabular}
\caption{The certified ladder. ``States'' is the in-edge-reachable boundary count
of the transfer operator. ``Tries'' is how many candidate numerators the search
tried before one passed; ``$\mu_H$ (float)'' is the value the power-iteration
search phase proposed, recorded for comparison and used nowhere in the proof. Receipts,
one line each with git revision, host, vector checksum and precision fields, are
in \repo{results/strip\_mu\_certificates.log}. $H=2$ is the independent anchor:
$\mu_2 = 1 + \sqrt2 = 2.41421356\dots$}
\label{tab:ladder}
\end{table}

The certified value is a floor. At $H = 2$ the certificate says $2.4142135$
where $\mu_2 = 2.4142136$ to seven places, and from $H = 14$ upward it falls
measurably short of the float value, reaching $6.543$ against $6.5465$ at
$H = 17$ for the reason given in \S\ref{sec:precision}. Clamping can only lower
the certified value; it can never falsify it.

\subsection{Soundness of the exact phase}
\label{sec:exact}

The tool has two phases and only the second enters the proof. Phase~1 is a
warm-started bisection on $x$ for $\rhoo(M(x)) = 1$, in double precision, using
the production column-sweep kernel; it proposes a candidate $x$ and a candidate
vector, and nothing in it enters the proof. Phase~2 scales the vector to integers
with the largest entry at $2^{\mathit{vbits}}$, sets
$\mathit{num} = \lfloor \mu_{\text{float}} \cdot \mathit{den}\rfloor$, and checks
$M(\mathit{den}/\mathit{num})\,v \ge v$ in \texttt{unsigned \_\_int128}.

Three things make that sound.

\paragraph{Every rounding is downward.} Each of the $H$ stages of processing a column
is a nonnegative linear map, hence monotone, and the ``cell placed'' branch
contributes $\lfloor \mathit{val}\cdot\mathit{den}/\mathit{num}\rfloor \le
\mathit{val}\cdot x$. So the computed vector is componentwise \emph{at most} the
true $M(x)v$, and a check that passes on the computed vector passes on the exact
product. The ``no cell'' branch has weight $1$ and is copied exactly.

\paragraph{Restriction is safe.} Only the finite support $S$ of the converged
eigenvector is held, and mass leaving $S$ is dropped. That certifies the
principal submatrix $M_S$, and $\rhoo(M_S) \le \rhoo(M)$, so a pass on $M_S$
implies a pass on $M$. Entries that quantise to zero are harmless for the same
reason: Collatz--Wielandt needs only $v \ge 0$ and $v \ne 0$, and every state is
still checked.

\paragraph{Overflow is loud.} Every multiply and add is guarded against $2^{126}$,
and a trip is a FAIL naming the state, never a silent wrap. The vector width is
auto-sized to $126 - \lceil \mathit{digits}\log_2 10\rceil - 6$ bits and a larger
explicit value is refused rather than clamped. Measured headroom through
$H \le 11$: the largest value ever held is $96.7$ bits against a $96$-bit vector
ceiling, so the merges add $0.7$ bits.

If the first candidate numerator fails, the tool searches for the largest one
that passes. The check is monotone in the numerator (a smaller numerator means
a larger $x$, and every entry of $M(x)$ is nondecreasing in $x$), so it brackets
by doubling and bisects, roughly $2\log_2(\text{gap})$ tries rather than
hundreds. The winning numerator is re-verified last, so the receipt describes the
rational that was actually certified. A failure is never absorbed into a
tolerance.

\subsection{The self-tests}
\label{sec:selftest}

Four self-tests run as a gate:

\begin{enumerate}
\item[\textbf{A}] $\mu_2 = 1+\sqrt2$, so $24142/10000$ must PASS. It does.
\item[\textbf{B}] $24143/10000$ exceeds $\mu_2$, so it must FAIL. It does, and
  the receipt names the offending state and prints $\mathit{lhs} < \mathit{rhs}$
  in full integers.
\item[\textbf{C}] One entry of a \emph{passing} vector is multiplied by $2^{30}$;
  the check must reject it. It does, naming the state.
\item[\textbf{D}] The all-zero vector satisfies $Av \ge v$ vacuously and would
  certify \emph{every} rational; it must be refused. It is.
\end{enumerate}

Test~D is the one that matters most, because it is the failure mode a reader
cannot see from a receipt. Bad arguments are likewise refused rather than
clamped: an out-of-range $H$, or a vector width above the arithmetic's headroom,
exits with the valid range printed.

\subsection{The precision budget}
\label{sec:precision}

The binding constraint at large $H$ is not time and not convergence. It is that
the eigenvector's dynamic range grows faster than the arithmetic ceiling.
Measured: the range spans $24.1$ bits at $H=8$, $52.5$ at $H=11$, $80.6$ at
$H=13$, $87.4$ at $H=14$, $97.8$ at $H=15$, $104.3$ at $H=16$ and $134.8$ at
$H=17$, against a vector ceiling of $96$ bits at seven decimal digits.

Entries whose true scale sits \emph{below} the quantisation floor are clamped to
$1$, and no few-bit increase reaches them --- verified by leaving the width at
$96$ and at $100$ and getting identical certified numerators, and again after
adding an eigenvector polish phase. The lever that works is the number of
decimal digits: a smaller denominator frees accumulator headroom for more vector
bits, since the cap is $126 - 6 - \log_2(10^{\mathit{digits}})$. Digits $6$, $4$
and $3$ buy widths $100$, $106$ and $110$, and at $H = 15$ and $H = 16$ that
lifts the whole eigenvector above the floor, so the certificate lands within one
unit in the last place of the float value.

At $H = 17$ it cannot: the range is $134.8$ bits, above the cap at \emph{every}
digits setting, so part of the tail is clamped and the certificate concedes about
$0.0035$ --- $6.543$ against a float $6.5465$. Harvesting that concession would
need at least $192$-bit accumulators. We did not pursue it, because each further
rung buys about $+0.05$ toward $\lambda \approx 7.11$, and $H = 18$ needs roughly
$35$~GB of transient build memory.

Forcing the width down measures the flooring loss directly. At $55$ bits at
$H = 10$, where the smallest entry becomes about $2^{10}$ instead of $2^{51}$,
the first candidate fails at four states with a minimum ratio of $0.99994$: the
binding state is short by one unit in the last place, so the flooring loss is
$O(1)$, not $O(H)$.
The bracket-and-bisect then lands on a slightly weaker rational in twenty tries,
at a wall time indistinguishable from the full-precision run.

\subsection{Scope}
\label{sec:honest}

What is certified is the inequality $\mu_H \ge \mathit{num}/\mathit{den}$,
\emph{given} that $M(x)$ is the height-$\le H$ strip transfer operator. The
operator itself is trusted, not certified here. It is the production kernel,
exercised by its own gates, and the resulting $\mu_H$ are cross-checked for
$H \le 11$ against roots of independently computed fixed-height generating
functions, agreeing to $9$--$10$ significant digits --- every digit those receipts
record. The certificates do not rest on that agreement; they are exact integer
arithmetic on the operator. That the operator is the strip operator is the one
thing here taken on trust.

Two engines were run against each other on every shared rung: an older hash-map
program and a frozen-successor-array kernel roughly $250\times$ faster reproduce
the receipts field for field at $H = 12, 13, 14$: same numerator, same state
count, same minimum ratio to all nine printed digits.

\subsection{An independent estimate of \texorpdfstring{$\lambda$}{lambda}}
\label{sec:estimate}

This subsection is measurement, not proof, and is used nowhere above.

A sliding three-point power-law fit $\mu_H = \lambda - cH^{-p}$ over the float
ladder gives $\lambda$ estimates that march monotonically downward as the center
rises: $9.54, 8.41, 7.92, 7.66, 7.50, 7.40, 7.33, 7.29$ for centres
$H = 5,\dots,12$, with the exponent $p$ tending to about $1$, a $1/H$
finite-size correction. The sequence converges toward $\approx 7.11$, the value the
differential-approximant extrapolation of the term ratios gives.

It is \emph{structurally} independent: strip spectra
and row-sum ratios share no data and no algorithm. It is not yet sharper than the
ratio estimate, still $7.29$ at $H = 12$, but it corroborates it.

Two caveats. Reaching $\mu_H$ through count ratios
$T(n,H)/T(n-1,H)$ instead of eigenvalues is badly under-converged for
$H \gtrsim 6$; at $n = 36$, $H = 18$ that ratio is $7.63$, above $\lambda$. And
extrapolation quality is uncertain: full Richardson in $1/H^2$ gives
$7.41$ while a local two-point fit gives $6.64$--$6.71$, so higher-order
corrections matter and the ladder is too short for a controlled estimate. The
rigorous \emph{lower bound} use needs no extrapolation and is unaffected by any
of this.

\section{The upper bound}
\label{sec:upper}

\subsection{Elementary counting: \texorpdfstring{$\lambda \le 12.2$}{lambda <= 12.2}}
\label{sec:crudebound}

\begin{proposition}
\label{prop:crude}
$\lambda \le 5^5/4^4 = 3125/256 \approx 12.2$.
\end{proposition}

\begin{proof}
Process opened cells FIFO from the scan-minimal root, so the animal is explored
by breadth-first search. A processed cell $u$ with parent $d$ is encoded by one
letter recording which cells of its \emph{frame} it newly opens, the frame being
$N(u)$ less $\{d\}$ less the shared set $N(u) \cap N(d)$. Shared sets have four
cells for an orthogonal $d$ and two for a diagonal one, so every frame has three
or five slots, and coverage holds because $N(u)$ is the parent together with the
shared set and the frame, with the shared set contained in $N(\mathrm{parent}(u))$
and the recursion terminating at the root, whose frame is all eight neighbors.
A replay decoder recovers the animal from the letter sequence, so each $n$-cell
animal maps injectively to $n$ letters of total weight $x^{n-1}y^n$. The letter
alphabet is the $5$-bit mask family with $\sum x^{\mathrm{opens}} = (1+x)^5$,
whence $a(n) \le [x^{n-1}](1+x)^{5n}\cdot O(1)$ and
$\lambda \le \min_b (1+b)^5/b = 5^5/4^4$ at $b = 1/4$.
\end{proof}

\begin{remark}[a repair, and what it cost]
\label{rem:repair}
An earlier version of this proposition ran the Redelmeier scan directly: include
or exclude each cell entering the untried frontier, note that an included cell
newly exposes only its four neighbors ahead in scan order, and conclude
$a(n) \le \binom{5n}{n}$. That argument is \textbf{false}, and the constant it
reaches is nevertheless right. ``Behind in scan order'' does not mean ``already
considered'': a re-entrant animal can contain a cell that is the ahead-neighbor
of no included cell, so it never enters the frontier at all and the map from
animals to decision strings is undefined on it. The witness is the hook
$\{(0,3),(0,2),(0,1),(0,0),(1,0),(2,0),(3,0),(3,1)\}$, in which $(3,1)$ is
reached only from below-behind. This is not a rare corner: the scheme misses
$2$ of $20$ animals at $n = 3$ and $96{,}065$ of $147{,}941$ at $n = 8$.

The BFS-frame derivation above proves the same constant soundly, and is
machine-verified --- round trip, distinctness, weight identity and alphabet
census over every animal with $n \le 8$, with a negative control that drops a slot.
\end{remark}

The bound is rigorous and needs no machinery. In the language of the twig
encodings below it is the ``Eden level'': the alphabet is every occupancy pattern
of the $K$ frame slots, so $A(x) = (1+x)^K$ and the bound is $K^K/(K-1)^{K-1}$.
The rook case $K = 3$ gives $6.75$ and the king case $K = 5$ gives $12.2$, which
is a useful cross-check on the derivation.

\subsection{Two shortcuts that do not work}
\label{sec:falsestarts}

Both are tempting, and one control kills both.

\paragraph{``$\lambda \le 8$ because each cell has eight neighbors.''} False.
The idea is to encode an animal by its spanning-tree parent directions, at most
eight per cell. The map is not injective: the parent-direction string does not
reconstruct the animal, because the spanning tree drops the animal's cycles.

\paragraph{``$\lambda \le 4^4/3^3 = 256/27 \approx 9.48$ from the growth tree.''}
Also false, and for the same reason. Each cell has at most four forward
king-neighbors, so the Redelmeier growth \emph{tree} is a $4$-ary tree on $n$
nodes; but animal $\mapsto$ tree is not injective, again because the tree drops
the exclude decisions on already-adjacent cells, which is exactly where the
cycles live. What is injective is the BFS-frame letter sequence of
Proposition~\ref{prop:crude}, and the difference between that and the naive
decision string is exactly the repair recorded above.

\paragraph{The control.} Both are caught by running the identical argument on
the square lattice, where it proves $\lambda_{\text{polyomino}} \le 4$ while
$\lambda_{\text{polyomino}} \approx 4.06$~\cite{barequetRoteShalah2016}. A king
bound has no published number to check against; its rook analogue does, so every
candidate bound here is run through that analogue first.

The would-be repair, a transfer matrix on the Redelmeier generation queue,
fails because that state is unbounded, its length growing with the perimeter.
There is therefore no simple finite-state valid bound between $12.2$ and the real
method: the Klarner--Rivest functional-equation
trick~\cite{klarnerRivest1973} exists to handle exactly that unbounded
dimension.

\subsection{The generic twig decomposition}

A generic twig decomposition was built and verified against brute force: the
enumerator reproduces \oeis{A006770}; the base fact
$a(n) \le G_8(n) \le n\,a(n)$ holds, where $G_8$ counts animals with a marked
cell whose W, SW, S, SE neighbors are forbidden, so that
$\lambda = \mathrm{growth}(G_8)$; the corner cell decomposes into twelve leaf
cases and three cut cases, and every leaf reduction and every cut convolution is
a numerically verified valid over-count.

And the bound is worse than doing nothing:

\begin{table}[H]
\centering
\begin{tabular}{l r r}
\toprule
decomposition & types & bound \\
\midrule
generic twig, window 1 & 28 & 12.50 \\
generic twig, window 1 + required-cell tracking & 28 & 13.00 \\
generic twig, window 2 & 8740 & 12.5382 \\
\bottomrule
\end{tabular}
\caption{Measured: the window is not the lever for a generic decomposition.
Compare Proposition~\ref{prop:crude}'s $12.2$.}
\label{tab:generic}
\end{table}

The window closes fast ($8740$ types in about a second), so the looseness comes
from the \emph{case routing} rather than from resolution: multi-neighbor cases
were routed as lossy linear re-marks instead of as convolutions, and that throws
away one of the two pieces every time.

\paragraph{The twig route is closed, at exactly the crude constant.} A later
round asked what the Klarner--Rivest twig mechanism gives on this lattice at
level $1$, with the pre-registered reading that a level-$1$ number above about
$10.5$ closes the route. It gives $5^5/4^4$ --- identical to
Proposition~\ref{prop:crude}, to the digit. The mechanism that lets twigs beat
the Eden level on the square lattice is deferral, and deferral is structurally
blocked under king adjacency: a deferred cell always lands in its child's shared
set. So the twig family cannot improve on the elementary bound here, and of the two
routes only \S\ref{sec:convolution}'s moves.

\subsection{The convolution system}
\label{sec:convolution}

The method that works is Bui's~\cite{bui2025}, ported to king adjacency. The
move is to case on \emph{one} free cell $d$ at a time:
\begin{equation}
  \varphi_T \;=\; \varphi_{T'} \;+\; \varphi_{T'}\cdot\varphi_{D},
  \label{eq:split}
\end{equation}
where $T' = T \cup \{d \text{ forbidden}\}$ and $D$ is the type that splits off at
$d$. The $d$-empty branch is an \emph{exact} partition, the animals of type $T$
with $d$ empty being exactly the animals of type $T'$, and the $d$-occupied
branch is a valid split over-count.

Two details make \eqref{eq:split} beat the generic version, and both
were the difference between $12.5$ and $9.3$:
\begin{enumerate}
\item multi-neighbor cases stay \emph{convolutions}, preserving both pieces,
  rather than becoming lossy re-marks;
\item the split-off type $D$ knows that all of the cut cell's other neighbors
  belong to the other side, so they are empty in the $d$-piece.
\end{enumerate}

All recurrences were verified as valid over-counts against brute force for split
windows $\le 2$; larger windows are valid by construction, since a larger window
only puts more genuinely-empty cells into the split type, which tightens the
bound while remaining an over-count.

\begin{table}[H]
\centering
\begin{tabular}{r r r}
\toprule
split window & types & $\lambda \le$ \\
\midrule
1 & 21   & 10.354 \\
2 & 185  & 9.402 \\
3 & 5930 & 9.306 \\
\bottomrule
\end{tabular}
\caption{The window lever saturates: window $2 \to 3$ moves the bound by
$0.096$.}
\label{tab:rd}
\end{table}

The other lever available in a decomposition of this kind, which free cell to
case on, has \emph{no effect} here. Lower-left, upper-right, diagonal
and orthogonal orderings all produce the identical $185$-type closure and the
identical $9.4022$ at window~$2$. The closure saturates to the same set
regardless of casing order, and the bound is a spectral property of the closed
system. (In Bui's richer multi-type systems this is a real lever; here it is
not.)

\subsection{The exact rational certificate}

\begin{theorem}
\label{thm:upper}
$\lambda \le 9.3154$, verified in exact rational arithmetic.
\end{theorem}

\begin{proof}[Method]
Let $F_x$ be the system map of the window-$3$ convolution system at the formal
variable value $x$; it is monotone on the nonnegative orthant, because every
recurrence is a sum of products of components with nonnegative coefficients.
Exhibit a rational vector $u > 0$ with
\[
  u_T \;\ge\; \bigl(F_x(u)\bigr)_T \qquad\text{for every type } T ,
\]
a \emph{super-solution}. By monotonicity $u$ dominates every iterate of $F_x$
from $0$, so the iteration is bounded at $x$, so the generating function has
radius of convergence at least $x$ and $\lambda \le 1/x$.

The certificate is $x = 2147/20000$, so $1/x = 20000/2147$, a $5930$-component rational vector with
denominators at most $10^6$. Every inequality was verified in exact
\texttt{Fraction} arithmetic, with no floating point anywhere in the chain, and
$1/x = 9.31532\dots$
\end{proof}

The gap between the certified $9.3154$ and the numerically located $9.306$ is a
deliberate $0.1\%$ safety margin, so that the super-solution has room; shrinking
it tightens the bound toward $9.306$ at the cost of larger fixpoint values. A
window-$2$ certificate, $x = 106251/10^6$, cross-checks at $\lambda \le 9.4117$
and also passes.

\paragraph{Formalization.} The certificate scheme is formalized in Lean~4:
\texttt{BuiSystem.certSum\_le}, \texttt{lambda\_le\_of\_buiSystem} and
\texttt{RatCert.lambda\_le} depend on standard axioms only, and the concrete
window-$3$ instance \texttt{lambda\_le\_of\_bui\_rd3} depends additionally on one
named \texttt{native\_decide} leaf, which is the certificate's arithmetic
validity check. The naming matters: a named leaf says exactly which finite
computation is being trusted, so the compiled-evaluation surface is enumerable
rather than diffuse. There is no \texttt{sorry} and no anonymous
\texttt{Lean.ofReduceBool} anywhere in the development.

\section{Where the remaining gap lives}
\label{sec:gap}

The bracket is $6.543 \le \lambda \le 9.3154$ around a value near $7.11$, so the
upper side is loose by about $31\%$. This section measures where that looseness
sits. Everything here is measurement against brute-force counts.

\paragraph{Diffuse, not concentrated.} Define
$\mathrm{slack}(T,n) = \mathrm{rhs}(T,n)/\mathrm{true}(T,n)$ for each recurrence
of the window-$2$ system, against brute-force counts. Over $152$ recurrences the
maximum slack is $1.222$, the median $1.144$ and the minimum $1.000$; nothing
exceeds $1.5$ and only four exceed $1.2$. The loosest are the far-reach split
types. There is no small set of bad types to hand-engineer: fixing the worst four
barely moves the bound.

\paragraph{It grows with $n$.} The median slack increases by about $+0.022$ per
term between $n = 7$ and $n = 8$. A purely \emph{local} over-count would be
constant in $n$; growth means the over-count \emph{accumulates with animal
size}. That is the distant-overlap signature: in $\varphi_{T'}\varphi_D$ the
$d$-piece and the other side may wrap around and collide arbitrarily far from
the cut, and no finite window can forbid it. The same wall saturated the window
lever in Table~\ref{tab:rd}.

\paragraph{The anchor is clean.} $G_8(n)/a(n) = 1.0, 1.0, 1.1, 1.23, 1.36, 1.50,
1.63$ for $n = 1,\dots,7$: growing, but only as fast as the polynomial prefactor
that $G_8 \le n\,a$ permits, so $\mathrm{growth}(G_8) = \lambda$ exactly. The
anchor leaks no exponential; all of the exponential looseness is in the
convolution.

\paragraph{What that rules out.} Multi-cell exact casing only \emph{relocates}
the over-count --- that is precisely the leaf/cut structure of the generic
decomposition, already measured at about $12.5$ in Table~\ref{tab:generic}.
Required-cell types can make cuts locally exact but cannot forbid distant
overlaps, so they meet the same wall; the best case is the demonstrated gap of
the polyomino method, about $14\%$, which would land near $8.1$ here. And the
king slack \emph{grows} with $n$ where the rook method's saturates, which says
eight-connectivity is genuinely worse, so $8.1$ is optimistic rather than
guaranteed.

\paragraph{An independent route, also closed.} Barequet, Ben-Shachar and
Osegueda's quasi-submultiplicativity~\cite{bbo2021} gives an upper bound from a
single enumerated term: with $a(40)$ known, a degree-$2$ correction polynomial
would yield $\lambda \le 7.745$, which would be a far better bound than ours.
(Degree $4$ already yields nothing.) The obstruction is that their lexicographic
split shatters a king comb into about $n/4$ components (measured), so the
reassembly code costs $n^{\Theta(n)}$, and the connected centroid split cannot
prescribe the two halves' sizes to within $O(1)$. This is the same connectivity
wall. A transfer check confirms the difficulty is not ours: the missing lemma
would also improve the published polyomino record from $4.5252$ to $4.3828$, so
it is known-hard rather than merely unattempted.

\section{Comparison}

\begin{table}[H]
\centering\small\setlength{\tabcolsep}{4pt}
\begin{tabular}{l l l l}
\toprule
bound & value & kind & source \\
\midrule
directed animals, king lattice & $\lambda \ge 3+2\sqrt2 \approx 5.828$ & closed form & \cite{bacher2015} \\
multi-directed animals, king lattice & $\lambda \ge 6.4752$ & numerical & \cite{bacher2015} \\
strip ladder, certified & $\lambda \ge \mathbf{6.543}$ & exact rational certificate & this paper \\
$\lambda$ (estimate) & $\approx 7.11$ & extrapolation & --- \\
convolution certificate & $\lambda \le \mathbf{9.3154}$ & exact rational certificate & this paper \\
elementary counting & $\lambda \le 12.2$ & closed form & this paper, Prop.~\ref{prop:crude} \\
\bottomrule
\end{tabular}
\caption{The state of rigorous knowledge about $\lambda$. Both new bounds are
finite objects a reader can re-check.}
\label{tab:compare}
\end{table}

For scale, the corresponding polyomino problem has been worked on for fifty
years: Klarner and Rivest's method~\cite{klarnerRivest1973} gives $4.65$, pushed
by Barequet and Shalah to $4.5252$, against a true value near $4.0626$ and a
proved lower bound above $4$~\cite{barequetRoteShalah2016}. Our $31\%$ upper gap
against their $11\%$ quantifies the eight-neighbor penalty.

\section{Open problems}

\begin{openproblem}[the real target]
Bring the upper bound below $8$, or below $7.5$. \S\ref{sec:gap} says the levers
inside this method class are exhausted; what is needed is a decomposition whose
over-count does not compound with $n$, which means one that can forbid distant
overlaps.
\end{openproblem}

\begin{openproblem}[certify further rungs]
$\mu_{18}$ and beyond need $\ge 192$-bit accumulators and roughly $35$~GB of
transient memory. Each rung buys about $+0.05$. The ladder converges to
$\lambda$, so in principle the lower side can be pushed arbitrarily close; in
practice the cost per rung grows about $3.5\times$.
\end{openproblem}

\section{Reproduction}
\label{sec:repro}

\begin{table}[H]
\centering\small
\begin{tabular}{@{}l R{0.56\textwidth}@{}}
\toprule
claim & check \\
\midrule
lower-bound certificates, Table~\ref{tab:ladder} & \texttt{make build/strip\_mu\_cert} then \texttt{./build/strip\_mu\_cert 2 17} \\
the four self-tests of \S\ref{sec:selftest} & \texttt{./build/strip\_mu\_cert --selftest}, or \texttt{make gate-strip-cert} \\
receipts & \repo{results/strip\_mu\_certificates.log} \\
$\mu_H$ against fixed-height GF roots & \repo{experiments/mu\_H\_precision\_audit.py} \\
the estimate of \S\ref{sec:estimate} & \repo{experiments/lambda\_from\_mu.py} \\
the false starts of \S\ref{sec:falsestarts} & \repo{docs/proofs/polyplet-upper-bound.md} \\
the generic decomposition, Table~\ref{tab:generic} & \repo{experiments/king\_bound.py}, \repo{experiments/king\_bound\_fast.py} \\
the convolution system, Table~\ref{tab:rd} & \repo{experiments/king\_bui.py} \\
the upper certificate, Theorem~\ref{thm:upper} & \repo{experiments/king\_certificate.py} \\
the slack audit, \S\ref{sec:gap} & \repo{experiments/king\_slack.py} \\
the concatenation route & \repo{results/growth-constant.md} \\
the Lean development & \texttt{cd polyplets \&\& lake build} \\
this paper's printed numbers & \repo{paper/verify\_l\_papers.py} \\
\bottomrule
\end{tabular}
\caption{Where each claim is checked. The two certificate checks are the only
ones that are part of the proof; everything else is a measurement or a
provenance record.}
\label{tab:repro}
\end{table}

The pipeline that produced the upper certificate was validated on two other
lattices before it was pointed at the king lattice: it reproduces the polyiamond
bound $\lambda_T \le 3.6108$ of Bui's companion paper, and it reproduces the
rook six-type bound $\lambda_2 \le 4.63$, including verifying Bui's own published
certificate $x = 100/463$ exactly. Only then was the king system built.

\bibliographystyle{plain}
\bibliography{shared/refs}

\end{document}
